\section{The actual global Paley--Wiener condensation and its half-plane range}
\label{sec:global-condensation}

The local quotient calculation does not determine the global orthogonal
projection.  This section computes more of that actual projection.  The
first result extends its fixed-cutoff convergence to every corrected
strip, using MA-WP.  The next two results identify the exact smooth
Paley--Wiener range and its Hilbert closure at polynomial weight zero.
The latter is a weighted version of the theta-image Hardy-space method,
with the source's particular range and centering retained.

Write
\[
 \mathscr S_{0,+}=\{f\in\mathscr S(\mathbb R):f(-t)=f(t),
                         \ f(0)=\widehat f(0)=0\},\qquad
 (\Theta f)(x)=2e^{x/2}\sum_{n\ge1}f(ne^x).
 \tag{MA-GP.1}
\]
The Fourier transform is the unitary transform with kernel
$e^{-2\pi it\xi}$.  MA-TR-101 identifies $\Theta$ with an actual
topological isomorphism onto $T_{1/2}\mathcal I$.
For $Y>0$ define
\[
 \begin{aligned}
 \mathscr S_Y&=\{f\in\mathscr S_{0,+}:
                       \operatorname{supp}\widehat f\subseteq[-Y,Y]\},\\
 \mathscr P_Y&=\Theta\mathscr S_Y,\qquad
 J_Y=\overline{\mathscr P_Y}^{\,H_{\alpha,s}},\qquad
 Q_Y^{\rm glob}=P_{J_Y}.
 \end{aligned}
 \tag{MA-GP.2}
\]
For $S=\{\infty\}$ these are the rational trivial-sector
Paley--Wiener levels of Masullo, Definition~7.3.2.1,
(7.3.2.1)--(7.3.2.7), p.~351, and Definition~7.3.2.3,
(7.3.2.12)--(7.3.2.17), pp.~354--355: the absolute and
invariant support conditions coincide in this one-place case.
The Hilbert projections are Definitions~11.2.1.1 and~11.2.1.3,
pp.~487--489, on the corrected actual theta image.  The notation
$Q_Y^{\rm glob}$ is deliberately different from the local additive
Fourier multiplier $Q_Y$ in MA-LA-014.

\paragraph{MA-GP-001 (actual fixed-cutoff condensation on every strip).}
For every $\alpha>0$ and $s\ge0$,
\[
 Q_Y^{\rm glob}\longrightarrow P_G\quad\hbox{strongly as }Y\to\infty.
 \tag{MA-GP.3}
\]
For $\chi,g\in C_c^\infty(\mathbb R)$ the actual operators satisfy
\[
 \bigl\|M_\chi C_gQ_Y^{\rm glob}
                  -M_\chi C_gP_G\bigr\|_1\longrightarrow0.
 \tag{MA-GP.4}
\]
For the fixed exhaustion $\chi_R(x)=\chi_0(x/R)$ of MA-WP-003,
with $g(0)=0$, the iterated trace limit is therefore
\[
 \lim_{R\to\infty}\lim_{Y\to\infty}
        \operatorname{Tr}(M_{\chi_R}C_gQ_Y^{\rm glob})
       =-\sum_\rho r_\rho\,\mathcal Mg(\rho-1/2).
 \tag{MA-GP.5}
\]
The retained multiplicities $r_\rho$ are exactly those of MA-WP.
For each fixed test $g$ one can also choose a diagonal
$Y=Y(R)\to\infty$ with the same limit.  This assertion gives no
prescribed quantitative relationship between $Y$ and $R$, and
does not identify this diagonal with the local coupled cutoffs.
No zero-location input is used.

\begin{proof}
Here is a direct density argument respecting both polar conditions.
Choose an even $\kappa\in C_c^\infty$ equal to one near zero.
For $f\in\mathscr S_{0,+}$ set
$f_N=\mathcal F^{-1}(\kappa(\,\cdot/N)\widehat f)$.
Then $f_N\to f$ in Schwartz space, $f_N$ is even and
$\widehat f_N(0)=0$, but $f_N(0)$ need not be zero.
Choose an even $v$ whose Fourier transform is compactly supported
away from zero and has integral one.  Thus
$v(0)=1$ and $\widehat v(0)=0$.
The functions
\[
                         \widetilde f_N=f_N-f_N(0)v
 \tag{MA-GP.6}
\]
have both polar conditions, are band limited, and converge to $f$
in Schwartz space.  Continuity of the actual theta map into
$\mathcal E$, followed by its continuous inclusion into
$H_{\alpha,s}$, proves that $\bigcup_YJ_Y$ is dense in $J$.
The spaces $J_Y$ increase.  For $j\in J$, approximate it by
$j_0\in J_{Y_0}$; then
$\|(1-Q_Y^{\rm glob})j\|\le\|j-j_0\|$ for $Y\ge Y_0$.
On $J^\perp$ every $Q_Y^{\rm glob}$ is zero.  This proves
(MA-GP.3).

MA-WP-003 and the actual ambient lift of MA-LIFT give
$T=M_\chi C_gP_G\in\mathcal S_1(H_{\alpha,s})$.
Since $P_GQ_Y^{\rm glob}=Q_Y^{\rm glob}$ and $TP_G=T$,
the difference in (MA-GP.4) is $T(Q_Y^{\rm glob}-P_G)$.
For a rank-one operator $u\otimes\ell$, its trace norm after
this composition is $\|u\|\,\|(Q_Y^{\rm glob}-P_G)^*\ell\|$,
which tends to zero by self-adjoint strong convergence.
Finite-rank approximation of $T$, with the uniform bound
$\|Q_Y^{\rm glob}-P_G\|\le1$, proves (MA-GP.4).
This is the usual trace-ideal continuity argument, now applied
to the actual all-weight operators rather than to an unproved
ambient lift.

Take traces and then apply MA-WP-003 to obtain (MA-GP.5).
For each $R\ge1$, (MA-GP.4) permits a choice $Y(R)\ge R$
for which the trace-norm difference is at most $1/R$.
This proves the stated diagonal version without a uniform
rate or a change in the order of the two limits.
\end{proof}

\paragraph{MA-GP-002 (the exact smooth support description).}
For $c=-\log Y$,
\[
 \mathscr P_Y
   =\{g\in T_{1/2}\mathcal I:
                              \operatorname{supp}g\subseteq[c,\infty)\}.
 \tag{MA-GP.7}
\]
This is an equality of actual smooth theta-image spaces.  It
does not assert that taking this intersection commutes with
Hilbert closure.

\begin{proof}
Put $u=\widehat f$.  The two polar conditions and Poisson
summation give the exact identity
\[
 (\Theta f)(x)=2e^{-x/2}\sum_{n\ge1}u(ne^{-x})
                              =(\Theta u)(-x).
 \tag{MA-GP.8}
\]
If $\operatorname{supp}u\subseteq[-Y,Y]$, every term vanishes
when $x< -\log Y$.  This proves one inclusion.
Conversely let $g=\Theta f$ have the support in (MA-GP.7).
Equation (MA-GP.8) gives $\Theta u(\log t)=0$ for $t>Y$.
The actual smooth M\"obius inverse of MA-TR-101 gives
\[
 2\sqrt t\,u(t)
    =\sum_{n\ge1}\mu(n)n^{-1/2}(\Theta u)(\log(nt)).
 \tag{MA-GP.9}
\]
The series converges with all derivatives on compact
positive intervals.  If $t>Y$, all its arguments satisfy
$nt>Y$, so it is zero.  Evenness gives $u=0$ outside
$[-Y,Y]$.  Thus $f\in\mathscr S_Y$.
The band-support condition is closed in Schwartz space;
MA-TR-101 also shows that $\mathscr P_Y$ is closed in
the smooth theta-image topology, as required by the
source's smooth closure convention.
\end{proof}

\paragraph{MA-GP-003 (the actual finite-level Hardy range).}
Assume $0<\alpha<1/2$ and $s=0$.  On the half-line $[c,\infty)$
first use the equivalent norm
\[
 K_c=L^2([c,\infty),e^{2\alpha x}\,dx),\qquad
 (U_cf)(y)=e^{\alpha(c+y)}f(c+y),\quad y\ge0.
 \tag{MA-GP.10}
\]
Let $\mathcal L$ denote the Laplace isometry onto
$H^2(\operatorname{Re}z>0)$, with boundary norm
$(2\pi)^{-1}\int|F(it)|^2dt$.
There is a pure Blaschke product $B_\alpha$ with zeros
\[
 z_\rho=\alpha+\tfrac12-\rho,
 \qquad \operatorname{Re}\rho<\tfrac12+\alpha,
 \qquad\hbox{multiplicity }m_\rho,
 \tag{MA-GP.11}
\]
such that the actual closed range satisfies
\[
            \mathcal L U_cJ_Y=B_\alpha H^2(\operatorname{Re}z>0).
 \tag{MA-GP.12}
\]
The equality concerns sets with equivalent Hilbert topologies;
the actual symmetric-weight orthogonal projection still needs
the Gram operator computed in the next section.  In particular,
finite Paley--Wiener levels impose all zero jets strictly to
the left of the upper strip edge, including any zeros to the
left of the full strip.  No hypothesis about whether such
exterior zeros occur is used.

\begin{proof}
\emph{The actual closure and shift invariance.}
On $[c,\infty)$ the ratio of the original weight to the weight
in (MA-GP.10) is
\[
 R_c(x)=e^{2\alpha|x|-2\alpha x},\qquad
            1\le R_c(x)\le e^{4\alpha\max(-c,0)}.
 \tag{MA-GP.13}
\]
Thus the two closures of $\mathscr P_Y$ are the same set.
By (MA-GP.7), $\mathscr P_Y$ is invariant under right
translations $g(x)\mapsto g(x-t)$ for $t\ge0$.
The map $U_c$ converts these to $e^{\alpha t}$ times
the unilateral translations on $L^2(0,\infty)$.

For clarity the Laplace normalization and surjectivity can
be checked directly: its reproducing kernel is
$K(z,\zeta)=\int_0^\infty e^{-(z+\overline\zeta)y}dy
=(z+\overline\zeta)^{-1}$.
With $z=(1+w)/(1-w)$, the weighted composition
\[
 F(z)\longmapsto \frac{\sqrt2}{1-w}
                         F\!\left(\frac{1+w}{1-w}\right)
 \tag{MA-GP.14}
\]
has reproducing kernel $(1-w\overline\omega)^{-1}$.
Density of the spans of the evaluation kernels identifies
it isometrically onto disk $H^2$; Plancherel gives the
stated half-plane boundary norm.

Write $\mathcal M_Y$ for the Laplace image of the actual
closure.  Invariance under unilateral translations makes
it invariant under multiplication by $e^{-tz}$, $t\ge0$.
For each vector their strongly continuous Bochner average with
weight $e^{-t}$ converges in Hilbert norm and defines the bounded
operator of multiplication by $(1+z)^{-1}$.
It is consequently invariant under
$(z-1)/(z+1)=1-2/(1+z)$.
Under (MA-GP.14) this is disk multiplication by $w$.
Beurling's invariant-subspace theorem therefore gives
$\mathcal M_Y=\Theta_Y H^2$ for an inner function $\Theta_Y$.
The space is nonzero by any nonzero source test and the
injectivity of MA-TR-101.

\emph{The exact common zero divisor.}
The relation between Laplace and Mellin transforms is
\[
 (\mathcal L U_cg)(z)=e^{zc}\mathcal Mg(\alpha-z).
 \tag{MA-GP.15}
\]
All core functions therefore vanish with the multiplicities
in (MA-GP.11), by the full smooth ideal description.
These are all the common zeros, with no extra multiplicity.
To verify that assertion using actual source tests, take
an even $u\in C_c^\infty((-Y,Y)\setminus\{0\})$ with
$\int u=0$, and put $f=\mathcal F^{-1}u$.
It belongs to $\mathscr S_Y$.  For
$s_z=z+1/2-\alpha$, termwise Mellin integration first in
$\operatorname{Re}s_z>1$, followed by analytic continuation,
gives
\[
 (\mathcal L U_c\Theta f)(z)
       =2e^{zc}\zeta(s_z)
                    \int_0^Y u(t)t^{s_z-1}\,dt.
 \tag{MA-GP.16}
\]
Here $\operatorname{Re}s_z>1/2-\alpha>0$ throughout the
half-plane; the zero of the integral at $s_z=1$ removes
the simple pole of zeta.
For any prescribed $s_z\ne1$ the integral can be made
nonzero subject to $\int_0^Y u=0$: the densities
$1$ and $t^{s_z-1}$ are linearly independent on any
positive open interval.  At $s_z=1$ one can instead
make $\int u(t)\log t\,dt\ne0$, so the removable value
of the product is nonzero.
Thus its common zeros are exactly the nontrivial zeros
of $\zeta(s_z)$ in this half-plane, to their exact orders.
The functional equation identifies that multiset with
(MA-GP.11), through $\rho\mapsto1-\rho$.
This uses the functional equation and the classical
containing strip, not RH or any narrower zero location.

\emph{No singular inner factor or positive delay.}
Every $g\in\mathscr P_Y$ is smooth, flat at $c$, and
decreases faster than all exponentials as $x\to\infty$.
Consequently $v=U_cg$, extended by zero to the negative
half-line, is smooth, flat at zero, and has all derivatives
in $L^1(0,\infty)$.
Integration by parts gives, for all $N$,
\[
 |(\mathcal Lv)(z)|\le C_N(1+|z|)^{-N}
                       \quad(\operatorname{Re}z\ge0).
 \tag{MA-GP.17}
\]
Its Laplace transform is entire by the same exponential
decay.  Therefore its disk image in (MA-GP.14) is bounded
and extends holomorphically across every circle arc
except the point $w=1$ corresponding to infinity.

Fix one nonzero core function.  The singular measure of
the common inner factor is bounded above by that of this
bounded disk function.  The elementary Poisson-kernel
argument in MA-HY-001 shows that it is supported on the
isolated boundary zeros of that function and at $w=1$:
on an arc with a nonzero analytic continuation,
boundedness of its Poisson integral makes the singular
measure absolutely continuous and hence zero there.
At an isolated zero of finite order, a positive atom
would force decay $\exp(-C/(1-r))$ along a radius,
contradicting the finite-order analytic lower bound.
Thus only an atom at $w=1$ can remain.  In half-plane
coordinates its factor is $e^{-dz}$ for some $d\ge0$.

Such a common factor would force every $U_cg$ to vanish
on $[0,d)$.  Indeed multiplication by $e^{-dz}$ under
the Laplace isometry is precisely the unilateral shift
by $d$, as is seen first by direct integration and then
by $L^2$ continuity and uniqueness.
There is a core function with no such delay.
Choose $u$ even and smooth, supported in $[-Y,Y]$, positive
on $(Y-\delta,Y)$ and flat at $Y$, with a compensating
negative bump below $Y/2$ so that $\int u=0$.
It vanishes near zero.  For $x>c$ sufficiently close to
$c$, (MA-GP.8) has only its $n=1$ term, and
\[
                  (\Theta\mathcal F^{-1}u)(x)
                           =2e^{-x/2}u(e^{-x})>0.
 \tag{MA-GP.18}
\]
Hence $d=0$.
The common inner factor is the pure Blaschke product
with the zeros already identified.  Beurling's theorem
now proves (MA-GP.12), including Blaschke summability.
\end{proof}

\paragraph{What the finite-level result does and does not identify.}
The source's support condition has now become an actual Hardy
range, independently of zero location.  For a left-exterior
point $\lambda$ with $\operatorname{Re}\lambda< -\alpha$,
evaluation on the supported half-line is continuous, but its
squared dual norm contains
\[
 \int_c^0e^{2(\operatorname{Re}\lambda+\alpha)x}\,dx,
 \tag{MA-GP.19}
\]
which diverges as $c\to-\infty$.  At equality it is $-c$.
Thus the extra finite-level constraints can disappear when
the union is completed.  This is consistent with both
(MA-GP.3) and MA-WP's actual full-strip multiplicities.
It supplies no license to omit those constraints from
a finite-level kernel or to replace its weighted projection
by the local additive Fourier multiplier.

\paragraph{Primary-source and Grothendieck placement.}
The density argument in Masullo, Proposition~7.3.2.2(1),
pp.~352--353, and the condensation statements
Theorems~11.2.1.2/11.2.1.4, pp.~488--489, were read at
their stated scope.  MA-GP-001 applies the repaired all-weight
ambient trace theorem; it does not rely on the source's
unrepaired nuclear lift.
The Hardy input is Beurling (1949), Factorization Lemma,
\S7, pp.~245--246, and Theorem~I, \S9, pp.~247--249,
already checked in MA-HY.  A closely related unweighted
theta-image closure is Burnol (2001; author manuscript 2000), \S2, the proof preceding
Theorem~2.6, author pp.~9--10; its adelic version is
Theorem~3.2, author pp.~14--15.  The present proof supplies
the shifted norm, the actual source support condition,
and the finite-level projection interface needed here.
No priority claim is made for the Hardy-space method.

In Grothendieck's nuclear framework, (MA-GP.4) is an actual
trace-ideal approximation before the Fredholm trace is taken,
as in the cote~1 comparison, scans~079--080 and~111--112,
ALPOGE\_BRIDGES BR-TR-01/02.  Those archive locators remain
inherited.  The finite-length modules of MA-WP still give
the spectral multiplicities; passing through an equivalent
half-line norm does not construct a new arithmetic motive
or an $\ell$-adic weight realization.

Burnol, J.-F. (2001). An adelic causality problem related
to abelian $L$-functions. \emph{Journal of Number Theory, 87}(2),
253--269. \url{https://doi.org/10.1006/jnth.2000.2616}.
Author manuscript (2000), version~3:
\url{https://arxiv.org/abs/math/0001013}.
